\begin{center}
    {\LARGE \textbf{LỜI GIẢI ĐỀ LUYỆN 05}}\\[6pt]
    {\large (Dành cho đội tuyển quốc gia)}
\end{center}

\section*{Bài 1. Hai đống}

Giả sử hai đống ta cần tạo là $P_1$ và $P_2$. Mỗi đống ban đầu rỗng.

Ta muốn mỗi đống có đúng \textbf{một giá trị lớn nhất}, gọi là:
\[
\max(P_1) = x,\quad \max(P_2) = y,
\]
và cần tối thiểu hoá $|x - y|$.

Do ta có thể hoán đổi tên hai đống, không mất tính tổng quát khi giả sử:
\[
x \ge y.
\]

Mục tiêu của ta:  
\[
\text{tìm } y < x \text{ gần } x \text{ nhất có thể}.
\]

\subsection*{1. Điều kiện khả thi của $x$}

Nếu tồn tại cặp $(A_i, B_i)$ sao cho:
\[
A_i > x \quad \text{và} \quad B_i > x,
\]
thì ta không thể chọn bất kỳ phần tử nào từ cặp này cho đống có giá trị cực đại là $x$.  
Do đó, $x$ chỉ hợp lệ nếu:
\[
x \ge \min(A_i, B_i)\ \ \text{với mọi $i$}.
\]

Gọi:
\[
M = \max_{i=1}^N \min(A_i, B_i).
\]

Vậy thì:
\[
x \ge M
\]
là điều kiện bắt buộc.

\subsection*{2. Hai cực đại không được thuộc cùng một cặp}

Ta không thể chọn $x$ và $y$ từ cùng một cặp $(A_i, B_i)$ vì mỗi cặp chỉ được phép chọn đúng một phần tử.  
Nếu $x$ và $y$ đến từ cùng một cặp, ta không thể đưa chúng vào hai đống khác nhau.

\subsection*{3. Chuẩn bị dữ liệu bằng cách sắp xếp}

Tạo mảng $S$ gồm $2N$ phần tử, chứa tất cả các giá trị $A_i$ và $B_i$.  
Đồng thời lưu lại mỗi phần tử thuộc về cặp nào.

Thí dụ:
\[
(1,2), (2,3), (1,4) 
\quad \Rightarrow \quad
S = [1,1,2,2,3,4].
\]

Sắp xếp $S$ theo giá trị.

\subsection*{4. Duyệt thử từng giá trị làm $x$}

Ta duyệt chỉ số $i$ từ $1$ đến $2N$ và thử đặt:
\[
x = S_i.
\]

Bỏ qua nếu $S_i < M$ vì không khả thi.

Khi $x = S_i$, giá trị tốt nhất cho $y$ là:
\[
y = S_{i-1},
\]
vì đó là phần tử gần $x$ nhất nhưng nhỏ hơn $x$.

\subsubsection*{Trường hợp xấu}

Nếu $S_{i}$ và $S_{i-1}$ thuộc cùng một cặp, ta không thể chọn cả hai làm cực đại của hai đống.  
Khi đó ta phải lùi lại một bước:
\[
y = S_{i-2}.
\]

\subsection*{5. Tính độ chênh lệch}

Mỗi lần xác định được $x$ và $y$, ghi nhận:
\[
\text{diff} = x - y.
\]

Đáp án cuối cùng là:
\[
\min \text{diff}.
\]

\subsection*{6. Tại sao mọi thứ đều hợp lệ khi $x \ge M$?}

Sau khi ta chọn được hai phần tử lớn nhất $x$ và $y$:

\begin{itemize}
    \item Đặt $x$ vào đống $P_1$.
    \item Đặt $y$ vào đống $P_2$.
\end{itemize}

Với mỗi cặp còn lại, luôn có ít nhất một phần tử:
\[
\le x
\]
(vì ta đã bảo đảm điều kiện $x \ge M$).

Ta đặt phần tử đó vào $P_1$.  
Như vậy:
\[
\max(P_1) = x, \quad \max(P_2) = y.
\]

Mọi ràng buộc đều được thoả mãn.

\subsection*{7. Độ phức tạp}

Sắp xếp mảng $S$ mất:
\[
O(N \log N).
\]

Duyệt toàn bộ $S$ và kiểm tra từng trường hợp mất:
\[
O(N).
\]

Tổng cộng:
\[
O(N \log N).
\]

\newpage

\section*{Bài 2. Tổ hợp thú vị}

\subsection*{Tiền đề kiến thức}
Bài toán sử dụng các kỹ thuật:
\begin{itemize}
    \item Tổ hợp (combinatorics),
    \item Chia kẹo vào hộp (stars and bars),
    \item Nguyên lý bù trừ (inclusion-exclusion).
\end{itemize}

\subsection*{1. Tính $F(A, D)$ cho một đa tập cố định}

Sắp xếp:
\[
A_1 \le A_2 \le \dots \le A_N.
\]

\subsubsection*{1.1. Mọi phép chia tối ưu đều là chia theo các đoạn liên tiếp}

\paragraph{Mệnh đề.} Tồn tại một phép chia tối ưu trong đó mỗi nhóm là một đoạn liên tiếp của dãy đã sắp xếp.

\paragraph{Chứng minh.}
Giả sử tồn tại $i < j < k$ sao cho $A_i$ và $A_k$ cùng nhóm, còn $A_j$ nằm ở nhóm khác.  
Chuyển $A_j$ vào nhóm của $A_i$ và $A_k$ sẽ:

\begin{itemize}
    \item Không tăng độ dài đoạn nên không tăng chi phí nhóm chứa $A_i, A_k$.
    \item Không tăng chi phí nhóm ban đầu của $A_j$ (thậm chí có thể giảm).
\end{itemize}

Lặp lại với mọi trường hợp tương tự ta thu được phép chia toàn bộ theo các đoạn liên tiếp.

\subsubsection*{1.2. Công thức chi phí theo hiệu hai phần tử liên tiếp}

Gọi phép chia tạo ra các đoạn $[L_t, R_t]$. Khi đó:
\[
\sum_t \left( D + A_{R_t} - A_{L_t} \right)
= kD + \sum_t \sum_{i=L_t}^{R_t-1} (A_{i+1}-A_i).
\]

Mỗi cặp liên tiếp $(A_i, A_{i+1})$ cho phép:
\[
\text{hoặc tách đoạn (tốn } D),\quad \text{hoặc nối liền (tốn } A_{i+1}-A_i).
\]

Vậy:
\[
F(A,D)
= D + \sum_{i=2}^N \min(A_i - A_{i-1},\, D).
\]

\subsection*{2. Biểu diễn đa tập bằng dãy hiệu}

Đặt:
\[
A_0 = 0,\quad A_{N+1} = M,\quad B_i = A_i - A_{i-1}.
\]

Đa tập $A$ được xác định duy nhất bởi:
\[
[B_1, B_2, \dots, B_{N+1}]
\]
thoả mãn:
\[
B_1 \ge 1,\qquad \sum B_i = M.
\]

Để thuận lợi ta thay $B_1$ bằng $B_1 - 1$, lúc này:
\[
B_i \ge 0,\quad \sum_{i=1}^{N+1} B_i = M - 1.
\]

Chi phí trở thành:
\[
F(A,D)
= D + \sum_{i=2}^N \min(B_i, D).
\]

\subsection*{3. Phân loại theo số lượng hiệu lớn}

Gọi $x$ là số lượng chỉ số $i$ (từ $2$ đến $N$) sao cho $B_i \ge D$.  
Chọn vị trí của $x$ phần tử này:
\[
\binom{N-1}{x}.
\]

Mỗi phần tử trong số đó có dạng:
\[
B_i = D + C_i,\quad C_i \ge 0.
\]

Các phần tử còn lại (giữa $2$ và $N$) phải thỏa:
\[
B_i < D.
\]

Chi phí:
\[
F(A,D) = D + xD + \sum_{\text{nhỏ}} B_i,
\]
nên gọi:
\[
K' = K - D,
\]
và yêu cầu:
\[
\sum_{\text{nhỏ}} B_i = K' - xD.
\]

\subsection*{4. Điều kiện tổng toàn bộ dãy}

Tổng các $B_i$ thỏa:
\[
B_1 + B_{N+1} + \sum_{i=2}^N B_i = M-1.
\]

Không mất tính tổng quát, ta giả sử $B_2,B_3,\ldots,B_{x+1} \ge D$ (các phần tử lớn).

Thay các $B_i$ đã phân loại vào:
\begin{align}
    &B_1 + B_{N+1} + (B_2 + \cdots + B_{x+1}) + \sum_{\text{nhỏ}} B_i = M-1 \\
    \Rightarrow &B_1 + B_{N+1} + \sum_{i=1}^x (D + C_i) + (K' - xD) = M-1 \\
    \Rightarrow &B_1 + B_{N+1} + \sum_{i=1}^x C_i = M - 1 - K'.
\end{align}

Số nghiệm của các biến tự do ($B_1, B_{N+1}, C_i$) là:
\[
\text{Stars and Bars}:\
\binom{(M-1-K') + x+2-1}{x+2-1}=\binom{M - K' + x}{x+1}.
\]

\subsection*{5. Xử lý phần ``nhỏ hơn $D$'' bằng Inclusion--Exclusion}

Phần $N-1-x$ biến $B_i$ còn lại phải thỏa:
\[
0 \le B_i < D,\quad \sum B_i = K' - xD.
\]

Nếu không giới hạn trên, số nghiệm là:
\[
\binom{(K' - xD) + (N-1-x) - 1}{(N-1-x)-1}.
\]

Dùng bao hàm--loại trừ cho điều kiện $B_i < D$:

Gọi $Q_j$ là số nghiệm trong đó có \emph{ít nhất $j$} biến vi phạm ($B_i \ge D$).  
Khi vi phạm $j$ biến:
\[
B_i = D + Z_i,\quad Z_i \ge 0.
\]

Khi đó số nghiệm:
\[
Q_j = \binom{N-1-x}{j}
      \cdot
      \binom{(K' - xD - jD) + (N-1-x)-1}{(N-1-x)-1}.
\]

Số nghiệm hợp lệ:
\[
P_0 = \sum_{j=0}^{N-1-x} (-1)^j 
\binom{N-1-x}{j} Q_j.
\]

\subsection*{6. Độ phức tạp}
Ta duyệt $x$ từ $0$ đến $N-1$, mỗi lần tính được trong $O(N)$.

Tổng độ phức tạp:
\[
\boxed{O(N^2)}
\]

\subsection*{Kết luận}
Bài toán được giải bằng:
\begin{enumerate}
    \item Biểu diễn đa tập qua hiệu liên tiếp.
    \item Công thức $F(A,D)$ dùng $\min$ của hiệu và $D$.
    \item Phân loại theo số hiệu $\ge D$.
    \item Dùng stars-and-bars cho nhóm không giới hạn và nhóm tự do.
    \item Dùng bao hàm--loại trừ để xử lý ràng buộc $B_i < D$.
\end{enumerate}

\newpage

\section*{Bài 3. Dãy XOR tăng chặt}

Ta muốn tìm độ dài lớn nhất của một \textbf{đoạn con tăng dần} sau khi áp dụng phép biến đổi:
\[
A_i \leftarrow A_i \oplus X \quad (\text{với cùng một } X \text{ cho mọi } i).
\]

Hãy xét một đoạn con bất kỳ:
\[
[A_L, A_{L+1}, \ldots, A_R].
\]
Để đoạn này tăng dần, ta chỉ cần các điều kiện:
\[
A_L < A_{L+1},\ A_{L+1} < A_{L+2},\ \ldots,\ A_{R-1} < A_R
\]
sau khi XOR với $X$.

Do đó, ta chỉ cần quan tâm đến \textbf{từng cặp kề nhau} $(A_i, A_{i+1})$.

\subsection*{1. Phân tích điều kiện $A_i \oplus X < A_{i+1} \oplus X$}

So sánh hai số nhị phân sau khi XOR chỉ phụ thuộc vào \textbf{bit cao nhất mà chúng khác nhau}.

Gọi $k$ là bit cao nhất sao cho:
\[
A_i \text{ và } A_{i+1} \text{ khác bit thứ } k.
\]

Ví dụ:
\[
4 = 100_2,\ 5 = 101_2 \quad\Rightarrow\ k = 0,
\]
\[
4 = 100_2,\ 7 = 111_2 \quad\Rightarrow\ k = 2.
\]

Khi so sánh hai số, bit cao nhất mà hai số khác nhau \textbf{quyết định số nào lớn hơn}.

\paragraph{Nếu ban đầu $A_i < A_{i+1}$}  
Thì ở bit $k$:
\[
A_i[k] = 0,\quad A_{i+1}[k] = 1.
\]
Để bất đẳng thức giữ nguyên sau XOR:
\[
X[k] = 0.
\]

\paragraph{Nếu ban đầu $A_i > A_{i+1}$}  
Thì ở bit $k$:
\[
A_i[k] = 1,\quad A_{i+1}[k] = 0.
\]
Để đảo dấu quan hệ cho đúng (để khiến $A_i \oplus X < A_{i+1} \oplus X$):
\[
X[k] = 1.
\]

\textbf{Các bit khác của $X$ không ảnh hưởng tới bất đẳng thức này.}

Như vậy, mỗi cặp kề nhau $(A_i, A_{i+1})$ đặt ra \textbf{một ràng buộc tại đúng một bit $k$ của $X$}.

\subsection*{2. Xây dựng mảng ràng buộc $B$}

Tạo mảng:
\[
B_i = \text{“ràng buộc của $X$ do cặp } (A_i, A_{i+1}).
\]

Ràng buộc có dạng:
\[
B_i = 
\begin{cases}
+(k) & \text{nếu } X[k] = 1 \text{ được yêu cầu},\\
-(k) & \text{nếu } X[k] = 0 \text{ được yêu cầu}.
\end{cases}
\]

\subsection*{3. Một đoạn con $[L, R]$ tăng dần khi và chỉ khi không có xung đột bit}

Xét đoạn con $[L, R]$ sau khi XOR.

Điều này tương đương với việc tất cả các ràng buộc:
\[
B_L,\ B_{L+1}, \ldots, B_{R-1}
\]
\textbf{đồng thời khả thi}.

Xung đột xảy ra khi:
\[
\text{có bit } k \text{ bị yêu cầu vừa phải } X[k]=0 \text{ vừa } X[k]=1.
\]

Nếu không có xung đột → đoạn con tăng dần.

\subsection*{4. Tìm đoạn dài nhất với hai con trỏ (two pointers)}

Giữ hai con trỏ $L$ và $R$, và một cấu trúc dữ liệu $S$ chứa các ràng buộc hiện tại.

Quy trình:

\begin{enumerate}
\item Đặt $L = 1$, $R = 1$.
\item Trong khi thêm $B_R$ \textbf{không gây xung đột} trong $S$, tăng $R$ lên 1 và thêm ràng buộc mới vào $S$.
\item Khi có xung đột, ta ghi nhận độ dài đoạn:
\[
\text{len} = R - L + 1,
\]
và cập nhật đáp án.
\item Tăng $L$ lên 1 và loại bỏ ràng buộc $B_L$ ra khỏi $S$.
\end{enumerate}

Điều quan trọng:
\[
\text{Điểm } R \text{ tối ưu cho mỗi } L \text{ là không giảm khi } L \text{ tăng.}
\]

Do đó thuật toán two-pointer là hợp lệ và chạy tuyến tính.

\subsection*{5. Độ phức tạp}

\[
O(N \log N)
\]
(vì multiset hoặc cấu trúc tương đương để quản lý ràng buộc).
